\documentclass[10pt,a4paper]{amsart}
\usepackage{amsmath,amssymb,amsthm,graphicx}
\usepackage{verbatim}
\newenvironment{reference}{\comment}{\endcomment}
%\newenvironment{reference}{}{}
\newenvironment{slogan}{\comment}{\endcomment}
\newenvironment{history}{\comment}{\endcomment}
% For commutative diagrams we use Xy-pic
\usepackage[all]{xy}

% We use 2cell for 2-commutative diagrams.
\xyoption{2cell}
% Theorem environments.
%
\theoremstyle{plain}
\newtheorem{theorem}[subsection]{Theorem}
\newtheorem{proposition}[subsection]{Proposition}
\newtheorem{lemma}[subsection]{Lemma}

\theoremstyle{definition}
\newtheorem{definition}[subsection]{Definition}
\newtheorem{example}[subsection]{Example}
\newtheorem{exercise}[subsection]{Exercise}
\newtheorem{situation}[subsection]{Situation}

\theoremstyle{remark}
\newtheorem{remark}[subsection]{Remark}
\newtheorem{remarks}[subsection]{Remarks}

\numberwithin{equation}{subsection}

% Macros
%
\def\lim{\mathop{\mathrm{lim}}\nolimits}
\def\colim{\mathop{\mathrm{colim}}\nolimits}
\def\Spec{\mathop{\mathrm{Spec}}}
\def\Hom{\mathop{\mathrm{Hom}}\nolimits}
\def\Ext{\mathop{\mathrm{Ext}}\nolimits}
\def\SheafHom{\mathop{\mathcal{H}\!\mathit{om}}\nolimits}
\def\SheafExt{\mathop{\mathcal{E}\!\mathit{xt}}\nolimits}
\def\Sch{\mathit{Sch}}
\def\Mor{\mathop{\mathrm{Mor}}\nolimits}
\def\Ob{\mathop{\mathrm{Ob}}\nolimits}
\def\Sh{\mathop{\mathit{Sh}}\nolimits}
\def\NL{\mathop{N\!L}\nolimits}
\def\CH{\mathop{\mathrm{CH}}\nolimits}
\def\proetale{{pro\text{-}\acute{e}tale}}
\def\etale{{\acute{e}tale}}
\def\QCoh{\mathit{QCoh}}
\def\Ker{\mathop{\mathrm{Ker}}}
\def\Im{\mathop{\mathrm{Im}}}
\def\Coker{\mathop{\mathrm{Coker}}}
\def\Coim{\mathop{\mathrm{Coim}}}

% Boxtimes
%
\DeclareMathSymbol{\boxtimes}{\mathbin}{AMSa}{"02}

%
% Macros for moduli stacks/spaces
%
\def\QCohstack{\mathcal{QC}\!\mathit{oh}}
\def\Cohstack{\mathcal{C}\!\mathit{oh}}
\def\Spacesstack{\mathcal{S}\!\mathit{paces}}
\def\Quotfunctor{\mathrm{Quot}}
\def\Hilbfunctor{\mathrm{Hilb}}
\def\Curvesstack{\mathcal{C}\!\mathit{urves}}
\def\Polarizedstack{\mathcal{P}\!\mathit{olarized}}
\def\Complexesstack{\mathcal{C}\!\mathit{omplexes}}
% \Pic is the operator that assigns to X its picard group, usage \Pic(X)
% \Picardstack_{X/B} denotes the Picard stack of X over B
% \Picardfunctor_{X/B} denotes the Picard functor of X over B
\def\Pic{\mathop{\mathrm{Pic}}\nolimits}
\def\Picardstack{\mathcal{P}\!\mathit{ic}}
\def\Picardfunctor{\mathrm{Pic}}
\def\Deformationcategory{\mathcal{D}\!\mathit{ef}}
\usepackage{cleveref}
\crefname{originallemma}{original lemma}{original lemmas}
\Crefname{originallemma}{Original lemma}{Original lemmas}
\makeatletter
\newlabel{lemma-additivity-invertible-subsheaf}{{02UH}{source}{}{sourcepointer-1}{}}
\newlabel{lemma-additivity-invertible-subsheaf@cref}{{[originallemma][0][]02UH}{[1][0][]source}{}{}{}}
\makeatother
\crefname{originallemma}{original lemma}{original lemmas}
\Crefname{originallemma}{Original lemma}{Original lemmas}
\makeatletter
\newlabel{lemma-bivariant-zero}{{02UC}{source}{}{sourcepointer-2}{}}
\newlabel{lemma-bivariant-zero@cref}{{[originallemma][0][]02UC}{[1][0][]source}{}{}{}}
\makeatother
\crefname{originallemma}{original lemma}{original lemmas}
\Crefname{originallemma}{Original lemma}{Original lemmas}
\makeatletter
\newlabel{lemma-flat-pullback-cap-cj}{{02U8}{source}{}{sourcepointer-3}{}}
\newlabel{lemma-flat-pullback-cap-cj@cref}{{[originallemma][0][]02U8}{[1][0][]source}{}{}{}}
\makeatother
\crefname{originalreference}{original reference}{original references}
\Crefname{originalreference}{Original reference}{Original references}
\makeatletter
\newlabel{situation-setup}{{02QL}{source}{}{sourcepointer-4}{}}
\newlabel{situation-setup@cref}{{[originalreference][0][]02QL}{[1][0][]source}{}{}{}}
\makeatother
\title{Multiplicativity of the total Chern class in a short exact sequence}
\author{The Stacks Project Authors}
\date{stacks-project:chow.tex:local-snapshot-20261012; GNU FDL 1.2 or later}
\begin{document}
\maketitle
\noindent\textit{Editorial note.} All delivered mathematical text is the author's own, copied from the pinned Stacks Project chapter file chow.tex; no statement or proof step was rewritten, reordered or removed.\par\medskip
\noindent\textit{Editorial note.} Preamble provenance: the pinned chapter file opens with \textbackslash{}input\{preamble\} and carries no preamble text of its own, so the project's single Stacks-wide preamble staged at sources/stacks-chow/preamble.tex is excerpted verbatim as three spans, lines 8-12 (the verbatim package together with the project's own reference, slogan and history display-environment definitions, including the author's commented-out older form of the reference environment), lines 14-18 (Xy-pic setup) and lines 154-217 (the native theorem declarations and mathematical macros). The lines 8-12 span is required because a counted statement or its surrounding text may use those display environments, which the project typesets as comments. Nothing was added, removed or redefined; the wrapper class and the amsmath/amssymb/amsthm/graphicx base packages are supplied by the assembler.\par\medskip
\noindent\textit{Editorial note.} Same-chapter pointers are located at their \textbackslash{}label lines inside the pinned chapter file; cross-chapter pointers are located in the project's own tag list staged in the snapshot (sources/stacks-chow/tags.txt), at the verbatim lines carrying their canonical tags.\par\medskip
\noindent\textit{Allowed standard prerequisites (editorial).}\begin{itemize}
\item Vanishing criterion for bivariant classes. Same chapter chow.tex, label lemma-bivariant-zero, invoked at proof line 7990. Reduces the identity of bivariant classes to the case of an integral base, where it suffices to check the identity after capping with the fundamental class.
\item The rank one case of additivity. Same chapter chow.tex, label lemma-additivity-invertible-subsheaf, invoked at proof line 7993. Provides the base case of the induction on the rank of the submodule, where the submodule is invertible.
\item The projective space bundle formula and flat pullback versus capping with a first Chern class. Same chapter chow.tex, labels lemma-flat-pullback-cap-cj (proof line 8018) and the projective space bundle formula used at proof line 8002. The first shows that the pullback of classes to the projective bundle of the submodule is injective, which is what allows the identity to be checked upstairs, and the second lets the cap of the fundamental class be identified with the pullback of the class of the base.
\end{itemize}\medskip
\ifdefined\hypertarget\hypertarget{sourcepointer-1}{}\fi
\noindent\textit{Source reference.} Original lemma 02UH: sources/stacks-chow/pinned-original-source.tex, source lines 7925--7925. Navigation pointer for \textbackslash{}ref\{lemma-additivity-invertible-subsheaf\} used inside the counted span (line 7993); the argument that consumes it is retained in the span.\par\medskip
\ifdefined\hypertarget\hypertarget{sourcepointer-2}{}\fi
\noindent\textit{Source reference.} Original lemma 02UC: sources/stacks-chow/pinned-original-source.tex, source lines 6585--6585. Navigation pointer for \textbackslash{}ref\{lemma-bivariant-zero\} used inside the counted span (line 7990); the argument that consumes it is retained in the span.\par\medskip
\ifdefined\hypertarget\hypertarget{sourcepointer-3}{}\fi
\noindent\textit{Source reference.} Original lemma 02U8: sources/stacks-chow/pinned-original-source.tex, source lines 7440--7440. Navigation pointer for \textbackslash{}ref\{lemma-flat-pullback-cap-cj\} used inside the counted span (line 8018); the argument that consumes it is retained in the span.\par\medskip
\ifdefined\hypertarget\hypertarget{sourcepointer-4}{}\fi
\noindent\textit{Source reference.} Original reference 02QL: sources/stacks-chow/pinned-original-source.tex, source lines 1372--1372. Navigation pointer for \textbackslash{}ref\{situation-setup\} used inside the counted span (line 7966); the argument that consumes it is retained in the span.\par\medskip
\begin{lemma}
\label{lemma-additivity-chern-classes}
Let $(S, \delta)$ be as in Situation \ref{situation-setup}.
Let $X$ be a scheme locally of finite type over $S$.
Suppose that ${\mathcal E}$ sits in an
exact sequence
$$
0
\to
{\mathcal E}_1
\to
{\mathcal E}
\to
{\mathcal E}_2
\to
0
$$
of finite locally free sheaves $\mathcal{E}_i$ of rank $r_i$.
The total Chern classes satisfy
$$
c({\mathcal E}) = c({\mathcal E}_1) c({\mathcal E}_2)
$$
in $A^*(X)$.
\end{lemma}

\begin{proof}
By Lemma \ref{lemma-bivariant-zero} we may assume that $X$ is integral
and we have to show the identity when capping against $[X]$.
By induction on $r_1$. The case $r_1 = 1$ is
Lemma \ref{lemma-additivity-invertible-subsheaf}.
Assume $r_1 > 1$. Let $(\pi : P \to X, \mathcal{O}_P(1))$
denote the projective space bundle associated to $\mathcal{E}_1$. Note that
\begin{enumerate}
\item $\pi^* : \CH_*(X) \to \CH_*(P)$ is injective, and
\item $\pi^*\mathcal{E}_1$ sits in a short exact sequence
$0 \to \mathcal{F} \to \pi^*\mathcal{E}_1 \to \mathcal{L} \to 0$
where $\mathcal{L}$ is invertible.
\end{enumerate}
The first assertion follows from the projective space bundle formula
and the second follows from the definition of a projective space bundle.
(In fact $\mathcal{L} = \mathcal{O}_P(1)$.)
Let $Q = \pi^*\mathcal{E}/\mathcal{F}$, which sits in an
exact sequence $0 \to \mathcal{L} \to Q \to \pi^*\mathcal{E}_2 \to 0$.
By induction we have
\begin{eqnarray*}
c(\pi^*\mathcal{E}) \cap [P]
& = &
c(\mathcal{F}) \cap c(\pi^*\mathcal{E}/\mathcal{F}) \cap [P] \\
& = &
c(\mathcal{F}) \cap c(\mathcal{L}) \cap c(\pi^*\mathcal{E}_2) \cap [P] \\
& = &
c(\pi^*\mathcal{E}_1) \cap c(\pi^*\mathcal{E}_2) \cap [P]
\end{eqnarray*}
Since $[P] = \pi^*[X]$ we
win by Lemma \ref{lemma-flat-pullback-cap-cj}.
\end{proof}

\end{document}
